\section{From electric correlations to microscopic normalization}
\label{sec:calibration}

The electric formula is accurate when marked points remain separated
by a mesoscopic distance. Counting an edge requires two opposite-color
marked points at distance one. The outward coupling bridges these
scales. We first obtain a pure-power amplitude for a neighboring pair,
then compare that pair with a disk. The latter comparison is the
normalization needed for occupation measures.

\subsection{A signed factor with positive local conditioning}

Fix the two neutral charge groups of
Proposition~\ref{prop:special-charges}. Each group has a core charge
$\alpha=\pi/2-\lambda$ and sufficiently many small negative
compensating charges. Place the compensators at distinct points
$u_j(n)$ on the axis by fixed, color-preserving lattice roundings
$u_j(n)=n\zeta_j+O(1)$, where the $\zeta_j$ are distinct fixed
locations away from the origin. The two cores lie within distance
$r=o(n)$ of the origin, in opposite colors. The first core is held
fixed. Write $e$ for the neighboring pair corresponding to a medial
segment, and write $r$ for the separated pair, with the appropriate
odd lattice separation. All rounding conventions preserve these
colors.

Let
\[
 B_* = \sum_{j\text{ remote}}\beta(t_j),\qquad
 \beta(t)=-A(t^2+2\lambda t),\qquad
 A=\frac{1}{2\pi(\pi-\lambda)}.
\]
Choose disjoint small patches of radius a fixed multiple of $n$ around
the remote points. Reweight only the loops wholly contained in these
patches, using their single-point weights
$v(t_j)=\cos(\lambda+t_j)/\cos\lambda$. These weights are positive.
Let $Z_*(n)$ be the mass of this local reweighting and let
$\Pp_n^*$ be its normalized probability law. The cutoff comparison
and the pure-nest asymptotic give
\begin{equation}\label{eq:proxy-mass}
 Z_*(n)\asymp\prod_{j\text{ remote}}F_n(v(t_j))
 \asymp n^{B_*}.
\end{equation}
Constants are uniform over the finitely many rounding classes used
below. We use the positive-winding expansion of the electric
correlation. Its remaining loop factor, denoted by $Y_r$, need not
be nonnegative.

Let $R$ be larger than the core separation by a fixed factor and
smaller than a sufficiently small multiple of $n$. The event $S_R(r)$
means that both core colors reach the tile rim of radius $R$.
The scalar weight of a loop enclosing just one core and no remote point
is zero, because $\cos(\lambda+\alpha)=0$. A completed loop inside
the $R$-disk enclosing both opposite-color cores contains another
loop separating them. Thus any nonzero configuration has $S_R(r)$.
Conversely, factoring out its indicator merely restricts the original
loop product. Consequently the full electric correlation satisfies
\begin{equation}\label{eq:core-factorization}
 C_r(n)=Z_*(n)\Pp_n^*(S_R(r))
             \E_n^*[Y_r\mid S_R(r)].
\end{equation}
The same identity holds for the neighboring pair $e$.

We record two uniform bounds used throughout this section. For every
fixed $j\ge1$,
\begin{align}
 \E_n^*[|Y_r|^j\mid S_R(r)]&\le C_j,
 \label{eq:macro-moments}\\
 \frac{\Pp_n^*(S_R(r))}{\Pp(S_R(r))}
 &=1+O((R/n)^{c_0}).
 \label{eq:proxy-local-probability}
\end{align}
Here $\Pp$ denotes the ordinary plane FK law; the same statements hold
with $r$ replaced by $e$. To justify the first bound, discard
configurations on which $Y_r=0$. Every non-unit factor remaining
on the other configurations belongs to a macroscopic loop.
A separating loop must enclose a remote point, and an ancestor
enclosing both cores contains such a separating loop. A loop charged
only at remote points but not included in the local reweighting has
left its remote patch. Each such loop therefore has a passage through
one of finitely many buffered regions of size comparable to $n$.
These buffers may be chosen away from the support of $S_R(r)$.
Even a loop with an arbitrarily large excursion has such a passage.
Conditional passage moments and the local-tilt estimate of
Input~\ref{inp:fk} control every fixed moment of the product.
Cutting the arm event back by a fixed factor leaves collars before
the passage buffers, so the estimates cost its ordinary probability,
not its reciprocal. This gives \eqref{eq:macro-moments} after
normalization. The relative collar comparison gives
\eqref{eq:proxy-local-probability}, because the nonnegative local
reweighting is supported at distance comparable to $n$.

\begin{lemma}\label{lem:signed-comparison}
Consider two opposite-color core pairs contained in a disk of radius
$s$, with the same remote charges and local reweighting. If
$Cs<R<cn$, then
\[
 \left|\E_n^*[Y_1\mid S_R(1)]-
             \E_n^*[Y_2\mid S_R(2)]\right|
 \le C(s/R)^{c_1}.
\]
The constants are independent of the microscopic separations of the
pairs.
\end{lemma}

\begin{proof}
Apply Proposition~\ref{prop:outward} to the two starts, with the common
nonnegative remote reweighting. On successful coupling the band and
the entire exterior agree. Exactly one global loop has the two
through-arcs in that band. It separates the opposite-color cores:
an inner closure separating them would contradict $S_R$. Its
enclosures of remote points are the same in both configurations.
Every other loop outside the band either encloses both cores or
neither, and the ancestors of the through-loop enclose both.
Inside the band no completed loop can enclose a core under $S_R$.
The factors $Y_1,Y_2$ therefore agree on successful coupling.
H\"older's inequality and \eqref{eq:macro-moments} bound their
contribution on coupling failure by a positive power of its probability.
Reducing the exponent in Proposition~\ref{prop:outward} proves the
claim. No positivity of $Y_i$ is used.
\end{proof}

\subsection{The neighboring-pair amplitude}

Set $x_1=-\beta(\alpha)$. The parameter identities are
\begin{equation}\label{eq:exponent-identity}
 2x_1+2A\alpha^2=x,
 \qquad
 x=\frac{\pi-2\lambda}{2(\pi-\lambda)}.
\end{equation}
For a sufficiently small fixed $\xi>0$, choose an odd separation
$r\asymp n^{1-\xi}$. Proposition~\ref{prop:special-charges} gives
\begin{equation}\label{eq:mesoscopic-core}
 \begin{split}
 C_r(n)
 &=M_* n^{B_*-x}r^{2A\alpha^2}(1+O(n^{-c_2}))\\
 &=M_* n^{B_*}r^{-2x_1}(r/n)^x(1+O(n^{-c_2})),
 \qquad M_*\ne0.
 \end{split}
\end{equation}
The amplitude depends on the fixed remote geometry, charges, and
rounding convention, but not on $n$ or $r$ in this regime. Indeed the
full Coulomb formula uses the actual pairwise distances; replacing
the remote distances from either core by their limiting values changes
its product by $O(r/n)+O(n^{-1})$. These errors can be included in
the power remainder. We will not use this equivalence with a fixed
$M_*$ when $r/n$ is held positive.

\begin{proposition}\label{prop:microscopic-amplitude}
For the fixed remote experiment there exists $c_*\ne0$ such that
\begin{equation}\label{eq:microscopic-amplitude}
 C_e(n)=c_*n^{B_*-x}(1+O(n^{-c_3})).
\end{equation}
In particular this limit holds through all positive integers $n$.
\end{proposition}

\begin{proof}
Take
\[
 r\asymp n^{1-\xi},\qquad R\asymp n^{1-b\xi},
 \qquad 0<b<1.
\]
All radii are chosen once for the comparison interval $[n,2n]$.
From \eqref{eq:core-factorization}, \eqref{eq:proxy-mass},
\eqref{eq:mesoscopic-core}, the one-arm bound, and
\eqref{eq:start-probabilities}, we obtain, for every small
$\delta>0$,
\begin{equation}\label{eq:expectation-lower}
 \left|\E_n^*[Y_r\mid S_R(r)]\right|
 \ge c_\delta(R/n)^x(r/R)^\delta.
\end{equation}
For example the upper estimate
$\Pp(S_R(r))\le C_\delta r^{-2x_1}(r/R)^{x-\delta}$
in the denominator gives precisely this inequality.
Its right side has order
$n^{-[bx+\delta(1-b)]\xi}$, whereas the signed comparison error
is $O(n^{-c_1(1-b)\xi})$. Choose $b>0$ sufficiently small and then
$\delta>0$ sufficiently small that
\[
 c_1(1-b)>bx+\delta(1-b).
\]
Lemma~\ref{lem:signed-comparison} is then a relative power-error
comparison of the two conditional expectations.

Apply the same argument with the remote geometry at any integer scale
$N\in[n,2n]$, using the \emph{same} core positions, $r$, and $R$.
Equation~\eqref{eq:proxy-local-probability} removes the remote tilt
from the arm probabilities with a power error. Uniformly in this
interval,
\begin{equation}\label{eq:fixed-radii-ratio}
 \frac{C_e(N)}{C_r(N)}
 =\frac{\Pp(S_R(e))}{\Pp(S_R(r))}(1+O(n^{-c_4})).
\end{equation}
The ordinary probability ratio on the right does not depend on $N$.
The mesoscopic formula therefore yields
\begin{equation}\label{eq:dyadic-amplitude-ratio}
 \frac{N^{x-B_*}C_e(N)}{n^{x-B_*}C_e(n)}
 =1+O(n^{-c_5}),\qquad n\le N\le2n.
\end{equation}
The denominator is nonzero for all sufficiently large $n$, since
\eqref{eq:expectation-lower} dominates the coupling error.

Write $a_n=n^{x-B_*}C_e(n)$. On a dyadic sequence, the consecutive
ratios of $a_n$ differ from one by a summable sequence. Their product
converges to a nonzero value, with tail error $O(n^{-c_5})$.
Equation~\eqref{eq:dyadic-amplitude-ratio} compares every intervening
integer to the preceding dyadic scale. Hence the entire sequence has
the same nonzero limit and the stated power error.
\end{proof}

\begin{proposition}\label{prop:sharp-arms}
For all radii above a fixed lattice scale,
\begin{equation}\label{eq:sharp-arms}
 \Pp(S_R(e))\asymp R^{-x},\qquad
 \pi_2(r,R)\asymp(r/R)^x.
\end{equation}
The estimates hold with fixed-factor changes of the inner and outer
rims and with a further bulk collar.
\end{proposition}

\begin{proof}
First choose $R=\theta n$, with $\theta>0$ a sufficiently small fixed
constant. Equations~\eqref{eq:core-factorization},
\eqref{eq:macro-moments}, and~\eqref{eq:microscopic-amplitude}
give the lower bound $\Pp(S_R(e))\ge c n^{-x}$.

For the reverse bound, compare the neighboring pair with an
opposite-color pair at separation $r\asymp aR$, where $a>0$ is a
small fixed constant. Use the full Coulomb formula, not the
$r=o(n)$ simplification with its limiting amplitude. Its magnitude
is uniformly comparable to
$n^{B_*}r^{-2x_1}(r/n)^x$ as first $n\to\infty$ and then
$a\downarrow0$. The nonzero amplitude and the remote-distance
factors vary continuously in this limit. It follows that
\[
 \left|\E_n^*[Y_r\mid S_R(r)]\right|
 \ge c_\delta\theta^x a^\delta.
\]
Choose $\delta<c_1$ and then $a$ small enough that
$Ca^{c_1}<\tfrac12 c_\delta\theta^x a^\delta$.
Lemma~\ref{lem:signed-comparison} gives a strictly positive lower
bound for the magnitude of the conditional expectation at $e$.
Equation~\eqref{eq:core-factorization} now gives
$\Pp(S_R(e))\le C n^{-x}$. Fixed-factor extensions cover all large
$R$, and finite energy covers the remaining bounded scales.

The first line of \eqref{eq:start-probabilities} identifies
$\Pp(S_R(e))$ with $\pi_2(C,R)$. Quasimultiplicativity therefore gives
\[
 \pi_2(r,R)\asymp
 \frac{\pi_2(C,R)}{\pi_2(C,r)}\asymp(r/R)^x.
\]
All comparisons used buffered events, so their collar and rounding
versions follow from the same fixed-factor estimates.
\end{proof}

\subsection{Calibration by disk probes}

A disk does not carry a prescribed electric charge. We therefore
introduce a diagnostic loop factor which agrees with the neighboring
pair's factor after successful outward coupling. Its purpose is only
to determine deterministic normalization constants.

Let $s=\epsilon n<R$. Under the same remote local reweighting, let
$S_R(s)$ be the disk two-arm event. If exactly one global loop touches
the disk and leaves the $R$-cut, assign that loop the factor it would
receive as a separator of the two cores, with its actual remote
enclosures. Give its ancestors the factors corresponding to enclosing
both cores, and give other loops no core charge. Divide out the local
remote weights already included in $\Pp_n^*$. Call the resulting
factor $Y_s$. In every other case set $Y_s=0$. In particular the
distinguished loop receives value zero when it encloses no remote
point. Define the deterministic quantity
\begin{equation}\label{eq:disk-diagnostic}
 W_s(n)=Z_*(n)\Pp_n^*(S_R(s))
                       \E_n^*[Y_s\mid S_R(s)].
\end{equation}
The moment bound \eqref{eq:macro-moments} also holds for $Y_s$:
every charged loop in a nonzero configuration spans between the disk
and a remote patch, or leaves a remote patch.

In the coupling of the disk start and the edge start, the factors
$Y_s$ and $Y_e$ coincide on success. The unique band through-loop
touches the disk, and no second loop can both touch the disk and leave
$R$. Replacing its inside portion preserves its remote enclosures;
the same is true for its ancestors. Thus
\begin{equation}\label{eq:disk-signed-comparison}
 \left|\E_n^*[Y_s\mid S_R(s)]-
             \E_n^*[Y_e\mid S_R(e)]\right|
 \le C(s/R)^{c_6}.
\end{equation}

\begin{lemma}\label{lem:disk-diagnostic}
For almost every fixed pair of scaled radii $0<\epsilon<\theta$
in the permissible range, with $s=\epsilon n$ and $R=\theta n$,
the limit
\begin{equation}\label{eq:diagnostic-limit}
 \lim_{n\to\infty}n^{-B_*}W_s(n)
\end{equation}
exists. Tile-rim roundings do not change the limit.
\end{lemma}

\begin{proof}
We use Corollary~\ref{cor:macroscopic-test}. First remove the nest weights
below radius $\delta n$ at each remote point, dividing by the
corresponding single-point nest partition functions. Their inward
coupling and cutoff insensitivity make the normalized error tend to
zero as $\delta\downarrow0$, uniformly in $n$.

There is a minor extra condition here: a loop deciding the disk
diagnostic might enter one of those small neighborhoods. It can be
discarded at a vanishing cost. Indeed the positive remote tilt produces
a surrounding screen before reaching the neighborhood with probability
tending to one. A loop arriving from the distant disk cannot cross
that screen. More annular trials reduce the exceptional probability,
and the moment bound for the remaining loop factors makes its weighted
contribution vanish. Thus, outside a vanishing exceptional event, the
diagnostic is determined by the truncated macroscopic configuration.

For fixed $\delta$, truncate excursions beyond a large fixed multiple
of $n$ and the number of charged macroscopic passages. The remaining
test uses finitely many loops of diameter bounded below after scaling.
It reads their ordered trajectories, their windings about the remote
points, and their intersections with the two disk rims. In a nonzero
diagnostic, ancestors of the distinguished loop can be recognized
among the remote-enclosing loops by whether they enclose the disk
center: any other such loop is disjoint from the disk, since otherwise
it too would touch the disk and leave $R$.

Input~\ref{inp:loops} gives convergence of these loop tests. Fixed
marked points are almost surely off all macroscopic loop traces.
For almost every choice of radii, no relevant loop is at a tangency
threshold, so intersection tests are continuity events. This follows
also by Fubini from the countable loop collection and its matching
topology. Bounded tile-rounding errors do not affect such tests.
The passage moment bounds remove the count and excursion truncations.
Finally let $\delta\downarrow0$. The pure-power asymptotics of the
remote nest factors identify their normalization with $n^{B_*}$,
proving \eqref{eq:diagnostic-limit}.
\end{proof}

\begin{proposition}\label{prop:disk-calibration}
There is a sequence $\epsilon\downarrow0$, a choice of radii
$R/n\asymp\epsilon^b$, and positive deterministic numbers
$L_\epsilon\asymp\epsilon^{-x}$ such that
\begin{equation}\label{eq:disk-calibration}
 \limsup_{n\to\infty}
 \left|
 \frac{n^x\Pp(S_R(e))}{L_\epsilon\Pp(S_R(\epsilon n))}-1
 \right|\le C\epsilon^{c_7}.
\end{equation}
The exponent $b>0$ may be chosen sufficiently small in advance to
satisfy any additional finitely many inequalities of the form
$(1-b)c>bx$, with $c>0$ fixed. The same constants apply to every
translate and rotation of the edge start, including color exchange.
\end{proposition}

\begin{proof}
Write $t=R/n$. Equations~\eqref{eq:core-factorization},
\eqref{eq:proxy-mass}, \eqref{eq:microscopic-amplitude}, and
\eqref{eq:sharp-arms} imply
\begin{equation}\label{eq:edge-expectation-size}
 \left|\E_n^*[Y_e\mid S_R(e)]\right|\asymp t^x.
\end{equation}
For $t\asymp\epsilon^b$, the relative error in
\eqref{eq:disk-signed-comparison} is bounded by
$C\epsilon^{(1-b)c_6-bx}$. Choose $b$ so that this exponent is
positive; it can simultaneously satisfy the stated finite collection
of further inequalities. The relative collar error in
\eqref{eq:proxy-local-probability} is $O(\epsilon^{bc_0})$.
Consequently
\begin{equation}\label{eq:disk-ratio-identity}
 \frac{n^x C_e(n)}{W_s(n)}
 =\frac{n^x\Pp(S_R(e))}{\Pp(S_R(s))}
                    (1+O(\epsilon^{c_7}))
\end{equation}
for a positive exponent $c_7$ after it is decreased if necessary.
Although the two correlations can be signed, their ratio in this
display is positive for small $\epsilon$, by the relative comparison.

Choose a sequence of continuity pairs for
Lemma~\ref{lem:disk-diagnostic}, with
$R/n\asymp\epsilon^b$. The sharp arm estimates show that the right
side of \eqref{eq:disk-ratio-identity} is comparable to
$\epsilon^{-x}$. Thus the limit of $n^{-B_*}W_s(n)$ is nonzero for
each sufficiently small chosen $\epsilon$. Define
\[
 L_\epsilon=
 \frac{c_*}{\displaystyle\lim_{n\to\infty}n^{-B_*}W_s(n)}.
\]
The numerator is the amplitude in
\eqref{eq:microscopic-amplitude}. This definition and
\eqref{eq:disk-ratio-identity} give positivity,
$L_\epsilon\asymp\epsilon^{-x}$, and
\eqref{eq:disk-calibration}.

All probabilities in \eqref{eq:disk-calibration} are ordinary plane
probabilities. Translation, square symmetries, and planar self-duality
identify them for the different edge locations and orientations, using
covariant tile-disk roundings. Therefore the calibration does not
depend on the checkerboard class or orientation of the probe.
\end{proof}
